\chapter*{Introduction générale}
\addcontentsline{toc}{chapter}{Introduction générale}
\markboth{Introduction générale}{}

Le commerce électronique s'est imposé au cours de la dernière décennie comme un canal de vente incontournable, et la pression concurrentielle qui en découle pousse les marchands à offrir une expérience client toujours plus réactive et personnalisée. Dans ce contexte, les assistants conversationnels fondés sur l'intelligence artificielle générative représentent une réponse naturelle : disponibilité continue, réponses instantanées, capacité à orienter le client dans un catalogue produits parfois volumineux. Toutefois, la mise en production de tels assistants soulève des difficultés qui dépassent largement la simple intégration d'un modèle de langage : isolement des données entre marchands (multi-tenance), fiabilité des réponses fournies (absence d'hallucination sur des informations factuelles comme le stock ou le prix), sécurité face aux tentatives de manipulation du modèle, et surtout, capacité pour l'exploitant de la plateforme à comprendre et piloter ce que fait réellement l'assistant.

C'est dans ce cadre que s'inscrit le présent stage de fin d'études, réalisé au sein de la société \textbf{Prodexo} sur une période de six mois, du 1\up{er} février au 1\up{er} août 2026. Le projet confié consistait à faire évoluer une plateforme d'assistant IA e-commerce existante -- initialement construite comme un socle technique fonctionnel mais partiellement instrumenté -- vers une version exploitable en conditions réelles par un administrateur de boutique, avec deux exigences transversales : que le comportement de l'assistant soit configurable sans intervention technique, et que les données présentées à l'administrateur (tableau de bord, statistiques, recommandations) soient systématiquement réelles et vérifiables, et non de simples valeurs d'exemple.

Ce rapport retrace la démarche suivie, depuis l'analyse du système existant jusqu'à la livraison d'une version consolidée du produit. Il s'organise en cinq chapitres. Le premier chapitre présente le cadre général du stage : l'entreprise d'accueil, le contexte et la problématique du projet, ses objectifs, ainsi que la méthodologie de gestion de projet adoptée. Le deuxième chapitre est consacré à l'analyse et à la spécification des besoins, avec l'identification des acteurs du système et de leurs besoins fonctionnels et non fonctionnels. Le troisième chapitre détaille la conception technique retenue : architecture générale, choix technologiques et modélisation. Le quatrième chapitre décrit la réalisation du projet, sprint par sprint, en mettant en avant les principales fonctionnalités développées. Le cinquième chapitre présente la démarche de tests et de validation mise en place. Le rapport se conclut par un bilan du travail accompli et les perspectives d'évolution envisageables pour la plateforme.
